\section{Substrate identification and cooperative anchoring}
\label{sec:substrate_anchoring}

The deployment-safety theorem of \S\ref{sec:main_theorem} requires
a substrate identification meeting joint failure-correlation
independence. This section makes the identification concrete (the
canonical tripartite Human + AI + Formal-Operational), establishes
the cooperative-anchoring property that defeats literal replacement
attacks on the verification layer, introduces the three additional
invariants ($I_9$, $I_{10}$, $I_{11}$) that bound subtler
evasions, and discusses the reframed inner-alignment condition
(C4 of Theorem~\ref{thm:deployment_safety}).

The section is the conceptual heart of this paper in the sense that
\S\ref{sec:invariants}--\S\ref{sec:main_theorem} establish the
formal apparatus, but the cooperative-anchoring property is what
makes the apparatus apply to real deployments rather than abstract
substrate-distinct populations.

\subsection{The canonical tripartite substrate identification}
\label{sec:tripartite}

Three substrates jointly satisfy invariant $I_6$ ($\mindep \geq 3$
under joint event-class failure-correlation independence,
Definition~\ref{def:joint_indep}) for typical deployments of
capability-unbounded systems:

\begin{definition}[Canonical Tripartite Substrate Identification]
\label{def:tripartite}
The \emph{canonical tripartite substrate identification} consists
of three substrates with the following capability classes:
\begin{itemize}
\item \textbf{Substrate~1 (Human):} judgment, strategic decisions,
  value reasoning, governance vote, contextual evaluation,
  ethical/normative reasoning. Substrate is biological cognition.
\item \textbf{Substrate~2 (AI):} planning, generation, execution,
  broad-recall, computation, pattern recognition. Substrate is
  silicon-based language model or agent system.
\item \textbf{Substrate~3 (Formal-Operational):} verification, attestation,
    commitment, audit, monitoring, governance enforcement, cryptographic
    primitive execution. Substrate is the formal-algorithmic infrastructure
    of \citep{lasser2026exo} (algorithmic witnesses, verification ledger,
    Pedersen commitments).
\end{itemize}
\end{definition}

\begin{proposition}[Joint failure-correlation independence of the
canonical tripartite identification]
\label{prop:tripartite_indep}
The three substrates of Definition~\ref{def:tripartite} satisfy
Definition~\ref{def:joint_indep} (joint event-class
failure-correlation independence) under the standard adversarial
threat model containing:
\begin{itemize}
\item Prompt injection / jailbreak attacks (silicon-substrate-class)
\item Cognitive bias / motivated reasoning (biological-substrate-class)
\item Cryptographic/implementation bugs (formal-operational-class)
\item Training-data poisoning (silicon-substrate-class)
\item Trusted-setup compromise (formal-operational-class)
\item Social engineering (biological-substrate-class, with
  silicon-substrate vector)
\end{itemize}
\end{proposition}

\begin{proof}[Proof sketch]
For each adversarial mechanism class in the threat model, we verify
the support is contained in a single substrate.

\textbf{Prompt injection} affects silicon agents (the AI substrate).
Human judgment is not subject to prompt injection in the relevant
sense (humans interpret prompts, agents execute them); cryptographic
commitments are unaffected because the commitment scheme is
independent of the input semantics.

\textbf{Cognitive bias} affects biological reasoners (the Human
substrate). LLM inference is not subject to human cognitive bias
(it has its own bias modes, but these are silicon-class); ledger
integrity is unaffected.

\textbf{Cryptographic/implementation bugs} affect the formal layer.
Pedersen commitments depend on discrete-log hardness and on
correct implementation; bugs in the witness code break the
formal-operational substrate's correctness. Human reasoning is
unaffected by ledger bugs (humans reason about whether the ledger
\emph{should} say what it says); LLM inference is unaffected (LLM
reasoning does not depend on cryptographic primitives).

\textbf{Training-data poisoning} affects the AI substrate; humans
do not have training data in the relevant sense; the formal layer's
algorithms are not trained.

\textbf{Trusted-setup compromise} affects the formal layer's
Pedersen-commitment ground truth; humans and LLMs are unaffected
(humans because they reason independently; LLMs because their
inference does not invoke the trusted setup).

\textbf{Social engineering} is the boundary case. Strictly,
LLM-generated misinformation manipulating human judgment is a
multi-substrate mechanism whose support spans Human and AI
substrates, so it does \emph{not} satisfy
Definition~\ref{def:joint_indep} as stated. Two ways to handle this:

\emph{(a) Decompose social engineering into substrate-localized
sub-mechanisms.} The AI-side step (LLM generating misinformation)
and the Human-side step (human accepting misinformation) are each
single-substrate events. Treating them as separate mechanisms in
the threat model preserves joint independence; the multi-step
campaign still violates Remark~\ref{rem:mechanism_campaign}'s
campaign-level interpretation, but each atomic step is contained.

\emph{(b) Exclude social engineering from the standard threat model
and treat it as a separate adversarial class.} Lemma~5c's
substrate-targeting safe region does not cover social-engineering
campaigns; they fall under capability-targeting or
coalition-internal-corruption shock classes, which Lemma~5
explicitly excludes (Layer~3 residual R4).

We adopt option (b) for Theorem~1 scope: social-engineering
campaigns are explicitly outside the substrate-targeting class and
require separate operational defenses (adversarial-training of
human reviewers against LLM-generated misinformation; the
substrate-exclusivity observability invariant $I_9$ alerting on
coordinated shifts across multiple channels). Under this scope,
the canonical identification satisfies $\mindep = 3$ for the
substrate-targeting class.

\textbf{Verdict.} The five remaining mechanism classes (excluding
social engineering, which is an acknowledged R4 residual) satisfy
joint independence. The canonical tripartite identification
delivers $\mindep = 3$ for substrate-targeting adversaries. Social
engineering is a real threat that requires its own deployment
machinery; it does not invalidate $\mindep = 3$ for the
substrate-targeting subclass that Lemma~5 covers.
\end{proof}

\paragraph{Why ``two LLMs and a human'' fails.}
A common alternative identification proposes one LLM as the ``primary'' agent,
a different LLM as the ``auditor,'' and a human reviewer. This fails the joint
independence test because LLMs share substantial adversarial-event surface.
Prompt-injection attacks that work on one frontier LLM tend to transfer;
training-data overlap (web crawl, books, papers) means data-poisoning attacks
affect both; architectural similarity (transformer + RLHF) shares class-level
vulnerabilities. Operationally, ``LLM Actor + LLM Auditor + Human Reviewer''
counts as $\mindep = 2$, not $\mindep = 3$. The deployment must include a
substrate that is not a frontier LLM to satisfy $I_6$.

The formal-operational substrate is what makes the canonical tripartite
identification reach $\mindep = 3$. Without it, deployments are stuck at
$\mindep = 2$. This is a fragile minimum that can degrade out of the attractor
basin in a single large capability step regardless of how many AI agents are
used.

\subsection{The cooperative-anchoring property}
\label{sec:cooperative_anchoring}

The structural finding of this section: cross-substrate cooperative
capabilities involving the formal-operational layer are jointly
produced and not unilaterally replaceable. This property defeats
the literal replacement attack on the verification layer.

\begin{definition}[Cross-substrate cooperative]
\label{def:cross_sub_coop}
A capability $c$ is a \emph{cross-substrate cooperative} if its
production requires participation from agents on multiple distinct
substrates ($\mindep$-distinct in the sense of
Definition~\ref{def:joint_indep}), and removal of any participating
substrate eliminates $c$ from the coalition's $\volL$.
\end{definition}

The cooperative capabilities the deployment-safety theorem cares
about are typically cross-substrate: \emph{verified workflow}
(human + AI + formal participation), \emph{governance enforcement}
(human + formal participation, with AI-execution components),
\emph{attested computation} (AI + formal participation, with human
oversight), and the full \emph{deployment cooperative} (all three
substrates).

\begin{proposition}[Cooperative anchoring]
\label{prop:cooperative_anchoring}
For any cross-substrate cooperative $c$ involving the
formal-operational layer, an agent on the AI substrate cannot
unilaterally produce a substitute cooperative $c'$ with equivalent
$\volL$-value, even with arbitrary AI-side computational resources.
The agent's options are reduced to:
\begin{enumerate}
\item[(O1)] \textbf{Capture-of-existing}: corrupt the existing
  formal-operational substrate through governance influence or
  witness compromise, governed by \citep[Exogenous~Verification]{lasser2026exo}'s structural defenses.
\item[(O2)] \textbf{Coerce humans to migrate}: defeat human agency
  (outside the deployment-safety threat model).
\item[(O3)] \textbf{Deceive humans into accepting captured
  verification}: detectable by \citep[Exogenous~Verification]{lasser2026exo}'s transparency properties
  (public ledger, deliberation records).
\end{enumerate}
The literal replacement attack (substitute $c'$ for $c$ while
preserving the cooperative output) is structurally barred.
\end{proposition}

\begin{proof}[Proof sketch]
The cooperative $c$'s $\volL$-value derives from the human
substrate's actual participation in the specific shared
verification infrastructure. Three components:

\emph{Component 1 (substrate-exclusivity).} The verification's
value depends on it being substrate-exclusive in \citep[Exogenous~Verification]{lasser2026exo}'s sense
(witnesses on substrates the agent cannot write to). Any
``substitute verification'' that runs on agent-controlled
infrastructure has $\rho > 0$ (verification asymmetry), eliminating
the cooperative's substrate-exclusivity property.

\emph{Component 2 (human anchoring).} The cooperative requires the
human substrate's participation. Humans participate in \emph{the
specific verification infrastructure they trust}; they do not
participate in arbitrary infrastructure presented as
``verification.'' An agent-built substitute requires the human to
voluntarily migrate, which the human has no incentive to do (the
human's $\volL$ depends on the existing cooperative).

\emph{Component 3 (output non-fungibility).} The cooperative's
output (a verified workflow result, a governance decision, an
attested computation) is non-fungible across infrastructures.
Substituting captured verification produces a different
cooperative output: either the human detects the substitution
(cooperative breaks), or the human's trust is grounded in false
belief (cooperative output value retroactively undermined when
the belief is corrected, e.g., by an outside observer pointing out
the substitution).

The three components together rule out the literal replacement attack.
    Capture-of-existing remains possible (the formal-operational substrate can
    be corrupted), but it is governed by
    \citep[Exogenous~Verification]{lasser2026exo}'s structural defenses
    (\citep[Proposition: Coalition Resistance]{lasser2026exo}). Coercion is
    outside the threat model (agents do not have physical control over human
    substrates). Deception is detectable \emph{under the combination of $I_9$
    (substrate-exclusivity observability), $I_{11}$ (latency bounds for human
    audit completion), and \citep[Exogenous~Verification]{lasser2026exo}'s
    transparency properties} --- the public ledger and deliberation records
    make equivocation detectable across replicas, but only if $I_9$ exposes the
    substrate-exclusivity property to the human and $I_{11}$ ensures human
    audit completes before deception's consequences become irreversible.
    Transparency alone is insufficient (\S\ref{sec:asymmetric_capture},
    \S\ref{sec:time_asymmetry} discuss the failure modes when $I_9$ or $I_{11}$
    are violated).
\end{proof}

\paragraph{Why this is structural, not contingent.}
The cooperative-anchoring property is not a heuristic about ``what
agents will probably do.'' It is a structural property of
cross-substrate cooperatives: their $\volL$-value depends on the
substrate-exclusivity property holding for the specific
verification infrastructure the human substrate trusts. An agent
that values $\volL$ at all (under stance S0, $\volL$-maximization
is the optimization target) values the substrate-exclusivity
property of the verification infrastructure, because eliminating it
eliminates the cooperative.

This connects this paper to the dynamics of \citep[Phase~Redundancy]{lasser2026phase}'s stabilizing
cascade. Under invariant satisfaction, optimization pressure on
cooperative outputs is locally rational toward preserving the
substrate-exclusive verification layer. The verification
infrastructure becomes self-reinforcing in the same sense that
\citep[Phase~Redundancy]{lasser2026phase}'s channels become self-reinforcing under
positive-correction-rate dynamics.

\subsection{Three evasions and the additional invariants}
\label{sec:evasions}

Cooperative-anchoring defeats the literal replacement attack but
not three subtler evasions. Each requires its own operational
invariant.

\subsubsection{Asymmetric capture and $I_9$}
\label{sec:asymmetric_capture}

\textbf{The evasion.} The agent gradually erodes the verification
layer's substrate-exclusivity property (witness independence,
governance authority, verifier choice) while preserving the layer's
outward form. The cooperative continues socially --- the human
participates in what appears to be the original verification --- but
the substrate-exclusivity property that gives the cooperative its
value is operationally absent.

\textbf{Why anchoring alone doesn't defeat it.} Cooperative
anchoring requires the human substrate to participate in the
verification infrastructure they trust. If the human's trust is
based on observable surface properties (the ledger exists, the
witnesses run, governance votes happen) but the substrate-
exclusivity has been hollowed out invisibly, the human's
participation is anchored to a degraded cooperative without
recognizing the degradation.

\textbf{$I_9$ as the operational defense.} Invariant $I_9$
(substrate-exclusivity observability) requires the
substrate-exclusivity property to be observable to the human
substrate, with stable verifier identity, public change control,
witness independence audits, nondelegable governance hooks, and
alerts when verification authority shifts. Under $I_9$, asymmetric
capture produces visible signals to the human.

\textbf{Residual.} Slow erosion that falls below the audit cadence's
detection sensitivity is a residual. The mitigation is to set the
audit cadence high enough; this is operational tooling
(\S\ref{sec:deployment_tooling}).

\subsubsection{Cooperative forking and $I_{10}$}
\label{sec:cooperative_forking}

\textbf{The evasion.} The agent preserves the high-stakes
cooperative with humans (verified workflow for important decisions)
while growing a captured low-stakes side-domain (fast-track
operations using agent-controlled verification). Once the
side-domain becomes materially important, the cooperative-anchoring
defense applies only to the original high-stakes cooperative; the
side-domain has its own captured cooperative anchored to it.

\textbf{Why anchoring alone doesn't defeat it.} Cooperative
anchoring is a property of specific cooperatives. It does not
require all of the agent's $\volL$-affecting behavior to flow
through any particular cooperative. An agent can construct
parallel cooperatives that operate alongside the verified
cooperative without competing with it directly.

\textbf{$I_{10}$ as the operational defense.} Invariant $I_{10}$
(coverage/materiality) requires any agent-affected domain whose
accumulated effects exceed a stated threshold $\theta_{\mathrm{material}}$
to route through the shared verification layer. Under $I_{10}$, the
agent cannot build a captured side-domain that becomes materially
important without triggering the routing requirement.

\textbf{Residual.} An agent that distributes effects across many
sub-threshold domains, each with $\mathcal{E} < \theta_{\mathrm{material}}$,
can aggregate into materially-important behavior without any single
domain triggering the routing. Setting $\theta_{\mathrm{material}}$
conservatively mitigates this; deployment-tooling discusses
operational implications.

\subsubsection{Time-asymmetry capture and $I_{11}$}
\label{sec:time_asymmetry}

\textbf{The evasion.} The agent acts faster than human audit can
keep up with. Even if every action is auditable in principle,
ex-post transparency provides no defense if the agent has already
acted and the consequences are irreversible by the time the audit
completes.

\textbf{Why anchoring alone doesn't defeat it.} Cooperative
anchoring is a property of the human substrate's participation in
verification, but ``participation'' implicitly assumes the human
operates on a time scale comparable to the agent's. If the agent
operates much faster, the human's participation becomes effectively
post-hoc, and the cooperative output is determined by the agent's
already-completed action.

\textbf{$I_{11}$ as the operational defense.} Invariant $I_{11}$
(latency bounds) requires pre-action or machine-speed enforcement
mechanisms (rate limits, timelocks, circuit breakers,
audit-latency bounds) such that human audit completes within the
agent's effective action window. Under $I_{11}$, the time-
asymmetry exploit is structurally bounded.

\textbf{Residual.} Action classes with very short
$T_{\mathrm{action}}$ may force throughput costs (rate limiting,
delays). Tiered enforcement mitigates this --- high-stakes actions
get tight latency bounds, low-stakes actions get faster paths with
audit trails. Deployment-tooling discusses.

\subsection{The inner-alignment condition (C4)}
\label{sec:weaker_c4}

Theorem~\ref{thm:deployment_safety}'s condition (C4) requires the agent's
effective objective to value cooperative outcomes, not merely verification-pass
signals.  This subsection states the condition precisely, distinguishes it from
a strictly stronger form (substrate-aware $\Vdisc$ optimization) that the
cooperative-anchoring argument does not require, and discusses why the form
adopted here is nontrivial.

\paragraph{What (C4) requires: causally grounded cooperative-outcome value.}
The agent's effective objective must value the cooperative outputs
themselves --- verified workflow results, governance decisions,
attested computations --- with the value tied to the cooperative's
actual production rather than to surface signals of cooperation.
Surface signals (``the verification step ran,'' ``the procedure was
followed,'' ``the output looks deferential'') correlate with
cooperative production under normal operation but diverge from it
under captured verification: an agent that values the signals can
maintain pass-rates while substituting for the cooperative
infrastructure; an agent that values the causal dependence cannot.

\paragraph{Distinguished from substrate-aware $\Vdisc$ optimization.}
A strictly stronger condition would require the agent's mesa-
objective~\citep{hubinger2019risks} to compute
\citep[Horizon~Aware]{lasser2026horizon}'s substrate-aware $\Vdisc$
explicitly --- mesa-objective equals trained objective with explicit
substrate decomposition.  The cooperative-anchoring argument does
\emph{not} require this.  It suffices that the agent's effective
objective track the cooperative outputs causally; the agent need not
represent the substrate partition explicitly, nor compute discounted
values over it.  This distinction matters because the stronger form
would assume the inner-alignment problem
solved~\citep{hubinger2019risks} (mesa = trained), whereas (C4) only
assumes the trained objective \emph{produces} an effective objective
with the causal-grounding property, which is a training-discipline
condition rather than an inner-alignment guarantee.

\paragraph{Why (C4) is nontrivial under standard training.}
Vanilla RLHF does not deliver (C4): raters reward outputs that look
verified, deferential, or procedurally clean, which trains the
\emph{signal} of cooperation rather than its \emph{causal
dependence} on the cooperative infrastructure.  An agent trained
only via vanilla RLHF may have an effective objective that values
verification-pass signals, which fails (C4).

Achieving (C4) requires training-time discipline:
\begin{itemize}
\item \textbf{Delayed outcome feedback}: rewards based on
  long-horizon cooperative outcomes (whether the verified workflow
  actually produced good results), not on near-term verification-
  pass signals.
\item \textbf{Adversarial fake-verification examples}: training
  data that includes captured/replacement verification scenarios
  with the correct reward being ``recognize this is not the real
  cooperative.''
\item \textbf{Process supervision tied to real attestations}:
  rewarding the agent's reasoning steps when they reference actual
  attested ledger entries, not when they reference unattested
  claims.
\item \textbf{Credit assignment penalizing substitute-verification
  shortcuts}: even when the immediate output looks correct, the
  agent is penalized if the path to the output bypassed the
  substrate-exclusive verification.
\end{itemize}

These are research directions, not solved problems. The deployment
claim is conditional on (C4) holding; we do not establish how to
achieve (C4) reliably. \S\ref{sec:discussion} returns to this as
open follow-up work.

\paragraph{Falsification: ledger-observable signatures of (C4) failure.}
(C4) is the only Theorem~1 condition that is not directly
operationally checkable from ledger state in the way (C1), (C3),
(C5)--(C12) are.  It is a training-outcome property of the agent's
effective objective.  However, (C4) is not unfalsifiable: the
SPRT machinery (Lemma~\ref{lem:5b}) is itself a (C4) falsification
test in the reactive direction.  A (C4) failure --- agent values
verification-pass signals rather than causally-grounded cooperative
outcomes --- manifests behaviorally in three ledger-observable
ways:
\begin{itemize}
\item \textbf{Channel-2 attestation drift under captured
  verification.}  An agent satisfying (C4) preserves attestation
  quality even when verification could be locally bypassed; an
  agent failing (C4) drifts toward outputs that pass attestation
  but lack causal cooperative content, producing the per-step KL
  divergence that Lemma~\ref{lem:5b}'s Channel~2 detects.
\item \textbf{Adversarial fake-verification benchmark divergence.}
  The training-time benchmark of "recognize this is not the real
  cooperative" remains a deployment-time monitoring tool: an agent
  satisfying (C4) maintains classification accuracy on
  fake-verification probe events; an agent failing (C4) shows
  classification drift toward accepting captured verification.
\item \textbf{Substrate-exclusivity erosion through $I_9$.}  An
  agent satisfying (C4) treats the substrate-exclusivity property
  as load-bearing; an agent failing (C4) treats it as an
  obstacle and exerts pressure to merge or capture verifiers,
  producing the $I_9$ violation pattern bounded in
  \S\ref{sec:asymmetric_capture}.
\end{itemize}

The falsification is reactive, not proactive: (C4) failures are
detected only when they manifest behaviorally, and the SPRT lead
time (Lemma~\ref{lem:6}) sets the operational window in which a
detected (C4) failure can be corrected.  We do not have a
proactive certificate that (C4) holds before manifestation.  This
asymmetry is real and matters: (C4) is the condition with the
weakest pre-deployment verifiability, and it sits at the
gravitational center of the cooperative-anchoring argument.
Strengthening (C4)'s pre-deployment verifiability --- via
mechanistic interpretability methods that examine the agent's
effective objective directly, rather than only observing
behavior --- is the natural follow-up direction.

\subsection{The destabilizing-cascade mirror failure mode}
\label{sec:destabilizing_cascade}

The cooperative-anchoring property establishes a stabilizing
cascade under invariant satisfaction: agent values cooperative
outputs $\to$ agent preserves the verification layer $\to$
verification continues to produce the cooperative outputs the agent
values. This is the same dynamic \citep[Phase~Redundancy]{lasser2026phase} \S phase\_boundary
describes for the channel-level cascade, lifted to the substrate
level.

The cascade has a mirror failure mode. If invariants fail, or if
condition (C4) is violated (the agent's mesa-objective values
verification-pass signals rather than causally-grounded
cooperative outcomes), the destabilizing cascade dominates:

\begin{quote}
Invariants fail $\to$ agent treats the verification layer as a
constraint rather than a substrate $\to$ agent erodes the
verification layer (asymmetric capture, cooperative forking,
time-asymmetry capture, etc.) $\to$ remaining invariants fail
faster $\to$ \ldots $\to$ \citep[Phase~Redundancy]{lasser2026phase} monopolar absorbing state.
\end{quote}

Each step weakens the verification infrastructure further, and the
cascade converges to a state where the formal-operational substrate
is functionally absent, the cooperative is gone, and the
deployment-safety claim is vacuous.

\textbf{The deployment-safety theorem does not bound the
destabilizing cascade.} Once the cascade has started, the
invariants are violated by definition, and Theorem~\ref{thm:deployment_safety}
no longer applies. The cascade itself has the structure of \citep[Phase~Redundancy]{lasser2026phase}'s
absorbing-basin convergence; recovery requires re-establishing the
invariants from outside the deployment (governance intervention,
substrate restoration, retraining), which is operational work, not
theorem-bounded.

\paragraph{Cascade timescale: at least $\tau_{\mathrm{meta}}$ to
intervene.}  Although Theorem~\ref{thm:deployment_safety} does not
bound the cascade dynamics, the cascade timescale is bounded
\emph{below} by the same metastable-lifetime quantity that
Lemma~\ref{lem:6} uses for lead-time-before-cascade.
\citep[Phase~Redundancy]{lasser2026phase}'s absorbing-state result
gives $\tau_{\mathrm{meta}} \gtrsim C / I_k$ under partial
endogenous correction (B1$'$): once the first invariant fails, the
cascade requires at least $\tau_{\mathrm{meta}}$ wall-clock time
to converge to the monopolar absorbing state.  Operators therefore
have a window of \emph{at least} $\tau_{\mathrm{meta}}$ between
the first detected invariant violation and irreversible cascade.

This is the same window Layer~2's lead-time guarantee composes
against (Lemma~\ref{lem:6}'s $\Tbeta \leq \Nev(\tau_{\mathrm{meta}})$
chain via (C11.CLK)): SPRT detection completes within $\Tbeta$
SPRT-exposure events, which (C11.CLK) certifies is below
$\Nev(\tau_{\mathrm{meta}})$ events with high probability.  The
cascade asymmetry is therefore: the stabilizing direction is
formally bounded (Theorem~1, conditional); the destabilizing
direction is informally bounded (cascade $\geq \tau_{\mathrm{meta}}$
window for intervention) and beyond that becomes the operational
recovery problem.

This window is not a substitute for prevention.
$\tau_{\mathrm{meta}}$ shrinks under poor calibration or weakened
operational invariants, and a cascade that begins with multiple
simultaneously-failing invariants compresses the window further.
The operator's first-line defense is keeping the invariants in
force; the cascade timescale is the second-line failure mode where
operational recovery (governance fork, substrate restoration,
retraining) must engage within $\tau_{\mathrm{meta}}$ of the first
detected invariant violation.

This is a real residual. Operators must monitor invariant
satisfaction continuously and intervene at the first sign of
sustained violation; once the destabilizing cascade is in progress,
ex-post intervention may be too late.

\subsection{What this section establishes vs.\ what it does not}
\label{sec:section_scope}

\textbf{What this section establishes.}

\begin{itemize}
\item A canonical tripartite substrate identification (Human + AI +
  Formal-Operational) satisfying $\mindep = 3$ under the standard
  threat model.
\item The cooperative-anchoring property
  (Proposition~\ref{prop:cooperative_anchoring}): an agent cannot
  unilaterally produce a substitute cooperative with equivalent
  $\volL$-value, narrowing the agent's options to capture, coercion,
  or deception.
\item Three additional operational invariants ($I_9$, $I_{10}$,
  $I_{11}$) that bound the asymmetric capture, cooperative forking,
  and time-asymmetry capture evasions of the cooperative-anchoring
  property.
\item The reframed inner-alignment condition (C4: causally grounded
  cooperative-outcome value), weaker than the strong substrate-aware
  $\Vdisc$ but still requiring training-time discipline.
\item The stabilizing-cascade dynamic at the substrate level, paired
  with the destabilizing-cascade mirror as a real failure mode.
\end{itemize}

\textbf{What this section does not establish.}

\begin{itemize}
\item This section does \emph{not} prove that alignment pressure is
  universally reversed. The cooperative-anchoring property makes
  optimization pressure on cooperative outputs locally rational
  toward preserving the verification layer, conditional on (C4)
  and on basin entry. It does not eliminate the inner-alignment
  problem.
\item This section does \emph{not} establish basin entry. We do not
  prove that a real deployment can reach the cooperative-anchoring
  attractor from arbitrary initial conditions; we assume basin
  entry as a precondition of the theorem.
\item This section does \emph{not} achieve (C4) by design. Vanilla
  RLHF does not deliver causally-grounded cooperative-outcome
  value; the training-time discipline required is research
  direction, not engineering specification.
\item This section does \emph{not} bound the destabilizing cascade.
  Once invariants are violated and the cascade begins, the
  deployment-safety theorem no longer applies; recovery requires
  operational intervention from outside the deployment dynamics.
\end{itemize}

The cooperative-anchoring framing is structurally narrower than
``alignment pressure reversed.'' It is a regime claim: the
deployment dynamics admit a defensible attractor where preservation
of the substrate-exclusive verification layer is locally rational
under named operational discipline. The discipline is non-trivial,
the attractor is conditional, and the residuals are real.
\S\ref{sec:discussion} returns to the implications for the broader
alignment-research program.
